\paragrafos{Ανάλογα ποσά: όταν το ένα διπλασιάζεται, διπλασιάζεται και το άλλο}{Α.6.3, «Ανάλογα ποσά – Ιδιότητες αναλόγων ποσών»}

\begin{anagke}
Στο παζάρι ο παππούς πουλάει το μέλι του: $8$ ευρώ το κιλό. Ένας πελάτης παίρνει ένα κιλό και πληρώνει $8$ €. Άλλος παίρνει δύο κιλά και πληρώνει $16$ €. Η Φανή κρατάει το βιβλίο και παρατηρεί ότι «όσο μεγαλώνουν τα κιλά, μεγαλώνουν και τα λεφτά, και μάλιστα πάντα το ίδιο: τα διπλά κιλά, διπλά λεφτά, τα τριπλά κιλά, τριπλά λεφτά». Η Πολυξένη το γράφει σε έναν πίνακα και βρίσκει ότι η τιμή διά τα κιλά κάνει πάντα $8$.

Ο Στέλιος ενθουσιάζεται και προτείνει: «Αυτό ισχύει παντού! Ο Νίκος είναι $12$ ετών και $1{,}55$ μ. Αν διπλασιάσουμε την ηλικία, στα $24$ θα έχει διπλάσιο ύψος, $3{,}10$ μ.». Όλοι γελούν και ο Νίκος φωνάζει: «Δεν είμαι καμηλοπάρδαλη!». Εδώ η Πολυξένη πάει στο τετράδιο: «Δύο μεγέθη που αλλάζουν μαζί δεν είναι αυτόματα ανάλογα. Πρέπει να ελέγχουμε».

Η παρέα θα μάθει σήμερα να ξεχωρίζει τα ποσά που μεταβάλλονται με τον ίδιο τρόπο (τα ανάλογα) από εκείνα που απλώς μεγαλώνουν μαζί.
\end{anagke}

\begin{lexi}{ανάλογα, συντελεστής αναλογίας}{ανάλογα, συντελεστής αναλογίας}
Τα \lexiword{ανάλογα} ποσά είναι αυτά που αλλάζουν «κατά τον ίδιο λόγο» (ανά + λόγος). Ο \lexiword{συντελεστής αναλογίας} είναι ο σταθερός λόγος $\dfrac{y}{x}$ των δύο ποσών· «συν-τελεστής» είναι αυτός που «τελεί» μαζί με το $x$, δηλαδή που το πολλαπλασιάζει.
\end{lexi}

\begin{metafrasi}
Κιλά μέλι και κόστος, $8$ € το κιλό. & $y=8\cdot x$\\
Διπλά κιλά, διπλά λεφτά. & $x\to2x\ \Rightarrow\ y\to2y$\\
Ο λόγος των ποσών είναι σταθερός. & $\dfrac{y}{x}=\alpha$\\
Το $12\%$ ενός ποσού $x$. & $y=0{,}12\cdot x$ (ανάλογα με $\alpha=0{,}12$)\\
\end{metafrasi}

\tmima{Η θεωρία}

\begin{orismos}{ανάλογα ποσά}
Δύο ποσά λέγονται \textbf{ανάλογα} όταν μεταβάλλονται με τέτοιον τρόπο ώστε, όταν οι τιμές του ενός πολλαπλασιάζονται με έναν αριθμό, και οι αντίστοιχες τιμές του άλλου να πολλαπλασιάζονται με τον ίδιο αριθμό.
\end{orismos}

Ισοδύναμα, δύο ποσά $x$ και $y$ είναι ανάλογα όταν οι αντίστοιχες τιμές τους δίνουν πάντα το ίδιο πηλίκο:
$$\dfrac{y}{x}=\alpha\quad\text{ή}\quad y=\alpha\cdot x,$$
όπου ο σταθερός αριθμός $\alpha$ (με $\alpha\neq0$) λέγεται \textbf{συντελεστής αναλογίας}. Στο μέλι ο $\alpha=8$ και $y=8x$.

\tmimasmall{Πώς ελέγχω αν δύο ποσά είναι ανάλογα}

\begin{enumerate}[leftmargin=1.6em,itemsep=0.3ex,topsep=2pt]
\item \textbf{Με τον ορισμό:} αν πολλαπλασιάζω το $x$ επί έναν αριθμό, το $y$ πολλαπλασιάζεται επί τον ίδιο;
\item \textbf{Με τον λόγο:} βρίσκω τους λόγους $\frac{y}{x}$ για κάθε ζεύγος του πίνακα. Αν είναι όλοι ίσοι, τα ποσά είναι ανάλογα.
\item \textbf{Με τον τύπο:} αν από το πρόβλημα έχω $y=\alpha x$, είναι ανάλογα.
\end{enumerate}

Αν δεν ισχύει έστω και μία φορά, τα ποσά δεν είναι ανάλογα. Επιπλέον, αν τα ποσά είναι ανάλογα και για $x=0$ έχουμε $y=0$.

\tmimasmall{Το ποσοστό είναι αναλογία}

Όταν το $y$ είναι ένα ποσοστό του $x$, δηλαδή $y=\frac{\alpha}{100}x$, τα δύο ποσά είναι ανάλογα με συντελεστή αναλογίας $\frac{\alpha}{100}$ ή $\alpha\%$. Έτσι το $12\%$ του $x$ είναι $y=0{,}12x$.

\tmimasmall{Ιδιότητες των αναλόγων ποσών}

Αν τα $x,\,y$ είναι ανάλογα, τότε:
\begin{itemize}[leftmargin=1.2em,itemsep=0.3ex,topsep=2pt]
\item Αν $(x_1,\,y_1)$ και $(x_2,\,y_2)$ είναι ζεύγη τιμών, τότε $\dfrac{y_1}{x_1}=\dfrac{y_2}{x_2}$, άρα έχουμε αναλογία και μπορούμε να βρούμε έναν άγνωστο όρο.
\item Το άθροισμα των τιμών ενός ποσού αντιστοιχεί στο άθροισμα των τιμών του άλλου: $(x_1+x_2,\ y_1+y_2)$ είναι κι αυτό ζεύγος τιμών.
\end{itemize}

\begin{prosoxi}
\begin{itemize}[leftmargin=1.2em,itemsep=0.3ex,topsep=2pt]
\item Δύο ποσά που μεγαλώνουν μαζί δεν είναι απαραίτητα ανάλογα: η ηλικία και το ύψος μεγαλώνουν μαζί, αλλά ο λόγος τους δεν μένει σταθερός.
\item Η σχέση $y=x+2$ δεν είναι αναλογία: για $x=0$ δίνει $y=2\neq0$. Στην αναλογία το $0$ πηγαίνει στο $0$.
\item Για να είναι ανάλογα, πρέπει να ισχύει για \emph{όλα} τα ζεύγη του πίνακα, όχι μόνο για τα δύο πρώτα.
\item Το εμβαδόν τετραγώνου δεν είναι ανάλογο της πλευράς του: πλευρά $1$ δίνει $1$, πλευρά $2$ δίνει $4$ (όχι $2$).
\end{itemize}
\end{prosoxi}

\tmima{Λυμένα παραδείγματα}

\paradeigma{Το μέλι του παππού}
\tixero Το μέλι κοστίζει $8$ € το κιλό.\par
\tithelo Να συμπληρώσω τον πίνακα για $0{,}5$, $1$, $2$, $3$, $5$ κιλά και να βρω πόσα κιλά αγοράζονται με $20$ €.\par
\possindeo Τα ποσά είναι ανάλογα με $y=8x$. Για $x=0{,}5$: $4$. $x=1$: $8$. $x=2$: $16$. $x=3$: $24$. $x=5$: $40$. Για $y=20$: $x=20:8=2{,}5$.\par
\apantisi $4,\,8,\,16,\,24,\,40$ € και $2{,}5$ κιλά.\par
\elegxos $2{,}5\cdot8=20$. Και $\frac{y}{x}=8$ σε κάθε ζεύγος.

\paradeigma{Είναι ανάλογα;}
\tixero Δύο πίνακες. (Α) $x$: $1,\,2,\,3,\,4$ και $y$: $2,\,4,\,6,\,8$. (Β) $x$: $1,\,2,\,3$ και $y$: $3,\,5,\,7$.\par
\tithelo Να δω ποιος περιγράφει ανάλογα ποσά.\par
\possindeo Ο λόγος $\frac{y}{x}$. Πίνακας Α: $\frac{2}{1}=\frac{4}{2}=\frac{6}{3}=\frac{8}{4}=2$, ίδιος. Πίνακας Β: $\frac{3}{1}=3$, $\frac{5}{2}=2{,}5$, $\frac{7}{3}\approx2{,}33$, δεν είναι ίσοι.\par
\apantisi Ο Α είναι ανάλογα ($\alpha=2$), ο Β όχι.\par
\elegxos Στον Β όταν το $x$ διπλασιάζεται από $1$ σε $2$, το $y$ γίνεται $5$, όχι το διπλάσιο του $3$.

\paradeigma{Βρίσκω το συντελεστή}
\tixero Τα $x,\,y$ είναι ανάλογα: για $x=4$ είναι $y=6$ και για $x=10$ είναι $y=15$.\par
\tithelo Ο συντελεστής και το $y$ για $x=14$.\par
\possindeo $\alpha=\frac{6}{4}=1{,}5$ και $\frac{15}{10}=1{,}5$. Για $x=14$: $y=1{,}5\cdot14=21$.\par
\apantisi $\alpha=1{,}5$ και $y=21$.\par
\elegxos $4+10=14$ και $6+15=21$ (ιδιότητα του αθροίσματος).

\paradeigma{Το διάλυμα}
\tixero Ένα διάλυμα έχει περιεκτικότητα $12\%$ σε αλάτι.\par
\tithelo Πόσο αλάτι υπάρχει σε $250$ g διαλύματος και πόσο διάλυμα περιέχει $9$ g αλάτι.\par
\possindeo Αλάτι $y=0{,}12x$, όπου $x$ το διάλυμα. Για $x=250$: $y=0{,}12\cdot250=30$ g. Για $y=9$: $x=9:0{,}12=75$ g.\par
\apantisi $30$ g αλάτι και $75$ g διαλύματος.\par
\elegxos $\frac{30}{250}=0{,}12$ και $\frac{9}{75}=0{,}12$.

\paradeigma{Ύψος και ηλικία}
\tixero Ο Νίκος στα $8$ ήταν $1{,}28$ μ. και στα $12$ είναι $1{,}55$ μ.\par
\tithelo Είναι ύψος και ηλικία ανάλογα;\par
\possindeo $\frac{1{,}28}{8}=0{,}16$ και $\frac{1{,}55}{12}\approx0{,}129$. Οι λόγοι δεν είναι ίσοι. Επίσης, η ηλικία αυξήθηκε $1{,}5$ φορές αλλά το ύψος μόνο $1{,}2$ φορές.\par
\apantisi Όχι, δεν είναι ανάλογα.\par
\elegxos Αν ήταν ανάλογα, στα $24$ ο Νίκος θα ήταν $3{,}1$ μ. Αδύνατο.

\tmima{Ασκήσεις}

\zoni{Το κατάλαβα}{πίνακες, συντελεστής, έλεγχος}

\ex Τα $x,\,y$ είναι ανάλογα με συντελεστή $\alpha=3$. Συμπλήρωσε: $x=0,\ 1,\ 2,\ 5,\ 0{,}4$ \ και βρες το $x$ όταν $y=9$ και όταν $y=4{,}5$.

\ex Τα $x,\,y$ είναι ανάλογα με $\alpha=\frac{2}{3}$. Βρες το $y$ για $x=0,\ 3,\ 6,\ 1$ και το $x$ για $y=10$.

\ex Εξέτασε αν τα $x,\,y$ είναι ανάλογα: α) $x$: $2,4,6,8$, $y$: $5,10,15,20$ \quad β) $x$: $1,2,3$, $y$: $2,3,4$ \quad γ) $x$: $3,6,9$, $y$: $1,2,3$ \quad δ) $x$: $1,2,4$, $y$: $5,10,21$.

\ex Βρες τον τύπο και πες αν τα ποσά είναι ανάλογα: α) το κόστος $y$ για $x$ τεμάχια των $1{,}5$ € \quad β) η περίμετρος $y$ τετραγώνου με πλευρά $x$ \quad γ) το εμβαδόν $y$ τετραγώνου με πλευρά $x$.

\ex Τα $x,\,y$ είναι ανάλογα με $x$: $5,\,8,\,12$ και $y$: $15,\,24,\,36$. Βρες τον συντελεστή.

\ex Σωστό ή λάθος; α) Αν διπλασιάσω το $x$ και το $y$ διπλασιάζεται, τα ποσά είναι ανάλογα. \quad β) Σε ανάλογα ποσά το $y$ είναι $0$ όταν το $x$ είναι $0$. \quad γ) Η $y=x+2$ είναι αναλογία. \quad δ) Το $y=0{,}2x$ είναι αναλογία με συντελεστή $20\%$. \quad ε) Το εμβαδόν ενός τετραγώνου είναι ανάλογο της πλευράς του.

\zoni{Το χρησιμοποιώ}{μέλι, βενζίνη, ύφασμα, ταχύτητα}

\ex \textbf{Το μέλι.} Συμπλήρωσε τα κόστη για $0{,}5,\ 1,\ 1{,}5,\ 2{,}5,\ 4$ κιλά μέλι των $8$ €. Πόσα κιλά αγοράζονται με $52$ €;

\ex \textbf{Η βενζίνη.} Ο θείος Λευτέρης αγοράζει βενζίνη προς $1{,}85$ € το λίτρο. Πόσο πληρώνει για $20$ και για $35$ λίτρα; Πόσα λίτρα παίρνει με $50$ €;

\ex \textbf{Το ύφασμα.} Η Αγγελική αγοράζει ύφασμα προς $4{,}8$ € το μέτρο. Πόσο κοστίζουν $2{,}5$ μ. και $0{,}75$ μ.; Πόσα μέτρα παίρνει με $33{,}6$ €;

\ex \textbf{Το διάλυμα.} Ένα διάλυμα έχει περιεκτικότητα $12\%$ σε αλάτι. Πόσο αλάτι υπάρχει σε $250$ g; Πόσο διάλυμα περιέχει $9$ g αλάτι;

\ex \textbf{Ο ΦΠΑ.} Ο ΦΠΑ είναι $24\%$ της αξίας. Βρες τον ΦΠΑ σε αξίες $50$ € και $200$ €. Για ποια αξία ο ΦΠΑ είναι $30$ €;

\ex \textbf{Ο ποδηλάτης.} Ο Νίκος ποδηλατεί με σταθερή ταχύτητα $18$ km την ώρα. Πόση απόσταση κάνει σε $0{,}5$, $2$ και $3{,}5$ ώρες; Σε πόσο χρόνο κάνει $27$ km;

\ex \textbf{Η μπογιά.} Ένα λίτρο μπογιάς καλύπτει $8$ m$^2$. Πόση επιφάνεια καλύπτουν $2{,}5$ λίτρα; Πόσα λίτρα χρειάζονται για $60$ m$^2$;

\zoni{Το σκέφτομαι}{ανάλογα ή όχι}

\ex Ο Στέλιος λέει ότι αφού ο Νίκος στα $12$ είναι $1{,}55$ μ., στα $24$ θα είναι $3{,}10$ μ. Γιατί σφάλλει;

\ex Ποια από τα παρακάτω ζεύγη ποσών είναι ανάλογα; α) τα κιλά και η τιμή \quad β) η ηλικία παιδιού και το βάρος του \quad γ) ο χρόνος και η απόσταση με σταθερή ταχύτητα \quad δ) η πλευρά και το εμβαδόν τετραγώνου \quad ε) η πλευρά και η περίμετρος τετραγώνου.

\ex Αν $y=2x$ και $z=5x$, δείξε ότι τα $y$ και $z$ είναι ανάλογα και βρες τον συντελεστή.

\ex Για το $y=3x$ τα $(2,\,6)$ και $(5,\,15)$ είναι ζεύγη τιμών. Έλεγξε ότι και το $(7,\,21)$ (το άθροισμα) είναι ζεύγος τιμών.

\zoni{Για το σπίτι}{ανάλογα στην κουζίνα και στο σούπερ μάρκετ}

\ex Βρες από τη συσκευασία ενός προϊόντος την τιμή ανά κιλό και φτιάξε πίνακα με τις τιμές για διάφορες ποσότητες.

\ex Βρες δύο ποσά της καθημερινής ζωής που είναι ανάλογα και δύο που μεγαλώνουν μαζί χωρίς να είναι ανάλογα.
